\section{Introduction}
\label{sec:intro}

Object detectors are increasingly used to automate visual inspection, where they must locate and classify several kinds of damage in the same image, such as cracks and potholes on roads or corrosion and cracking on bridges~\cite{shang2024automated,zheng2024deep,vision2025inspection}. Achieving reliable detection in these settings remains challenging, because the appearance of each damage type varies with illumination, surface texture, camera viewpoint and image quality, and a training set rarely covers this variation in full. The difficulty is compounded by differences among the categories themselves. Some are defined mainly by shape, such as thin longitudinal and transverse cracks~\cite{arya2024rdd}, whereas others are defined mainly by colour, such as rust and efflorescence on concrete~\cite{flotzinger2024dacl10k}. Consequently, improving detection in such settings requires exposing the detector to a wider range of visual conditions while preserving the cues that distinguish each category.

Data augmentation addresses the first part of this requirement by applying controlled edits, such as geometric, colour and image-quality transformations, to training images, so that the detector learns from more varied conditions without additional annotation. The effectiveness of an edit, however, depends on whether it preserves the cues that define each category. Colour jitter, for example, alters the very property that identifies rust while leaving the geometry of a crack intact, whereas Gaussian blur can erase a hairline crack while barely changing a pothole. The same edit can therefore provide useful variation for one category while removing essential information from another. In detection, this tension arises within a single training image rather than only between images: standard pipelines edit the whole image, so every object in it receives the same edit, and an edit that helps the detector learn one category can hinder another category in the same image. Such conflicts are easily overlooked, because augmentation is usually judged by accuracy averaged over classes, and gains in average accuracy can coexist with losses for individual classes~\cite{kirichenko2023}. This motivates augmentation that accounts for how each class responds to each edit, and that chooses the edit for an image according to the classes the image contains.

Existing augmentation methods have progressively moved toward this goal, but each stops short of it. AutoAugment and RandAugment search for or simplify one policy shared by the whole dataset~\cite{cubuk2019auto,cubuk2020rand}, and detection-specific extensions add edits restricted to annotated boxes or adapt the policy to object scale~\cite{zoph2020det,chen2021scale}; none of these policies, however, depends on the category of the objects being edited. Class-aware approaches recognize that categories may require different augmentations~\cite{tsuchiya2019,yoo2023,cheung2022adaaug}, and road-damage detection work has varied augmentation probabilities by damage type~\cite{rdd_detectron2}. Because these policies are fixed before or apart from detector training, they cannot follow changes in what each class needs as learning progresses. Online methods remove this restriction by adapting augmentation from the feedback of the evolving model, but the feedback they use introduces two further limitations. Tolerance-based controllers such as CTAugment and CUDA adjust augmentation according to whether the model still recognizes augmented examples~\cite{berthelot2020,ahn2023cuda}, which indicates what the model can already handle rather than whether training on the edit improves it. Benefit-based methods such as BDA, DDAS and MADAug measure that improvement through a held-out loss~\cite{lu2023bda,liu2021ddas,hou2023madaug}, but they compute it over all classes together, so an edit that helps frequent classes can be preferred even when it harms a rare one. Moreover, the class-aware methods among them are formulated for classification, where each training sample carries one label, and therefore do not specify which class should govern an image that contains several.

This gap matters in detection because the classes in one image can disagree. A detector may learn well from blurred images of potholes while the thin cracks in the same images lose the structure that identifies them. A policy guided by pooled feedback sees only the net effect of blur and may apply it to both, whereas a policy keyed to a single class must ignore the other. What is needed instead is a measure of how much each edit helps each class learn, combined into one decision for the image that respects every class it contains.

To meet this need, we propose Class-Utility Augmentation (CUA), an online augmentation policy that selects the edit for each training image from the classes present in it. CUA is built around the \emph{class utility} of an edit, defined as the estimated amount that the edit helps a given class learn. CUA estimates this utility during training with a first-order argument from training-influence analysis~\cite{pruthi2020tracin}: if a gradient step on the edited image would reduce a class's loss on a held-out \emph{reference set} more than a step on the original image would, the edit currently benefits that class. For each training image, CUA then averages the utilities of the classes present, removes any edit that is reliably harmful to one of them, and samples a single edit from the remainder, with ``no edit'' always available as a candidate of zero utility (Fig.~\ref{fig:teaser}). Because the utilities are re-estimated from the current detector, the preferred edits can change as training proceeds, and because CUA requires no policy network and no separate search phase, it adds only a periodic gradient computation to standard training.

The value of CUA rests on the premise that classes respond differently to the same edit, so we examine this premise before evaluating the method. Across road damage (RDD2022), bridge damage (dacl10k) and everyday objects (PASCAL VOC)~\cite{everingham2010voc}, and with two detector families, Faster R-CNN~\cite{ren2017faster} and RetinaNet~\cite{lin2017focal}, we train detectors with one edit at a time to measure its effect on each class, and with an edit applied only to the objects of one class to separate its direct effect on that class from its effect on the others. We then compare CUA with shared-utility augmentation, an otherwise identical policy that averages utility over all classes, so that any difference between the two reflects class information alone. \todo{One-sentence summary of findings: which edits affect classes in opposite directions; CUA minus shared-utility on test macro AP50 with 95\% intervals; rank agreement between estimated utilities and measured effects.} In summary, this paper makes the following contributions:

\begin{enumerate}
    \item Class-Utility Augmentation, an online augmentation policy for multi-class object detection that estimates how much each edit helps each class learn and selects one edit per image from the classes it contains, with a no-edit option and a per-class harm block.
    \item A controlled measurement of how individual edits affect each class in detection, across three domains and two detector families, which also separates the effect of an edit on the class it is applied to from its effect on the other classes in the image.
    \item A comparison with a shared-utility counterpart that differs from CUA only in class information, together with standard augmentation baselines, and a test of whether the estimated class utilities agree with the measured per-class effects.
\end{enumerate}

